% !TEX program = xelatex
\documentclass[aspectratio=169]{beamer}

\usetheme{Madrid}
\setbeamertemplate{footline}{}
\setbeamertemplate{navigation symbols}{}
\setbeamertemplate{caption}[numbered]

\usepackage{amsmath,amssymb}
\usepackage{graphicx}
\graphicspath{{media/}}
\usepackage{booktabs}

\setbeamercolor{takeaway}{fg=black,bg=structure.fg!8}

\newcommand{\mur}{\mu_{\mathrm r}}
\newcommand{\Tc}{T_{\mathrm c}}
\newcommand{\Tapp}{T_{1/2}^{\mathrm{app}}}
\newcommand{\takeaway}[1]{%
  \vspace{0.1em}
  \begin{center}
    \begin{beamercolorbox}[wd=0.94\textwidth,sep=0.35em,center,rounded=true]{takeaway}
      \textbf{#1}
    \end{beamercolorbox}
  \end{center}
}

\title[Ferromagnetism Lab]{Ferromagnetism and the Curie Transition}
\subtitle{Hysteresis loops as measurements of material response}
\author[I. Feldman, T. Bleher]{Itay Feldman \and Tom Bleher}
\institute[Tel Aviv University]{School of Physics and Astronomy\\Tel Aviv University}
\date{May 2026}
\titlegraphic{%
  \includegraphics[
    page=1,
    trim=185 650 185 55,
    clip,
    width=0.20\textwidth
  ]{/Users/tomble/Desktop/personal/data/Gemini/lab_tau-43e8004eb4267093.pdf}%
}

\begin{document}

\begin{frame}
  \titlepage
\end{frame}

\begin{frame}{Two Experiments, One Question}
  \large
  \begin{columns}[T,onlytextwidth]
    \begin{column}{0.48\textwidth}
      \textbf{Iron toroid}
      \vspace{0.5em}
      \begin{itemize}
        \item Map the peak envelope of $B(H)$.
        \item Extract an effective $\mur(H)$.
        \item Test the copper-gap route to $\mu_0$.
      \end{itemize}
    \end{column}
    \begin{column}{0.48\textwidth}
      \textbf{Monel sample}
      \vspace{0.5em}
      \begin{itemize}
        \item Record hysteresis loops while warming.
        \item Reduce each loop to a magnetization proxy.
        \item Locate the apparent Curie transition.
      \end{itemize}
    \end{column}
  \end{columns}
  \takeaway{The phenomena were clear; the final accuracy was set by systematic effects.}
\end{frame}

\begin{frame}{Physics In One Slide}
  \large
  \begin{columns}[T,onlytextwidth]
    \begin{column}{0.48\textwidth}
      \textbf{Ferromagnetic response}
      \vspace{0.5em}
      \[
        B=\mu_0(H+M), \qquad \mur=\frac{B}{\mu_0 H}
      \]
      \begin{itemize}
        \item Domains align with field.
        \item Hysteresis and saturation appear.
        \item $\mur$ is not constant.
      \end{itemize}
    \end{column}
    \begin{column}{0.48\textwidth}
      \textbf{Curie transition}
      \vspace{0.5em}
      \[
        M(T)\rightarrow0 \quad \text{near}\quad \Tc
      \]
      \begin{itemize}
        \item Thermal fluctuations destroy spontaneous order.
        \item Finite drive field smears the transition.
        \item We therefore report an operational midpoint.
      \end{itemize}
    \end{column}
  \end{columns}
\end{frame}

\begin{frame}{Measurement Logic}
  \small
  \begin{columns}[T,onlytextwidth]
    \begin{column}{0.48\textwidth}
      \centering
      \includegraphics[width=0.86\linewidth,height=0.42\textheight,keepaspectratio]{experiment.png}
      \vspace{0.2em}

      \textbf{Iron: absolute calibration}
      \[
        H=\frac{N_1}{2LR_x}\Delta V_x,\qquad
        B=\frac{R_yC}{2N_2A}\Delta V_y
      \]
    \end{column}
    \begin{column}{0.48\textwidth}
      \centering
      \includegraphics[width=0.86\linewidth,height=0.42\textheight,keepaspectratio]{experiment2.jpeg}
      \vspace{0.2em}

      \textbf{Curie: relative calibration}
      \[
        X\equiv V_x\propto H,\qquad
        M^*=c_YV_y-c_HV_x
      \]
    \end{column}
  \end{columns}
  \takeaway{$V_x$ measures drive field; the integrated pickup voltage measures magnetic response.}
\end{frame}

\begin{frame}[plain]
  \centering
  \includegraphics[width=0.95\paperwidth,height=0.94\paperheight,keepaspectratio]{fig_BH_mu.pdf}
\end{frame}

\begin{frame}{Iron: Magnetization Envelope}
  \centering
  \includegraphics[width=\textwidth,height=0.60\textheight,keepaspectratio]{fig_BH_mu.pdf}
  \vspace{0.1em}

  {\small
    $\mur^{\max}=2.18(8)\times10^3$ at
    $H=167(4)\,\mathrm{A/m}$, $B=0.458(12)\,\mathrm T$;
    final point $B=1.243(33)\,\mathrm T$.
  }
  \takeaway{Peak-to-peak data give a magnetization envelope, not a full hysteresis loop.}
\end{frame}

\begin{frame}[plain]
  \centering
  \includegraphics[width=0.95\paperwidth,height=0.94\paperheight,keepaspectratio]{fig_copper_gap.pdf}
\end{frame}

\begin{frame}{Iron: Copper-Gap Test of $\mu_0$}
  \centering
  \includegraphics[width=\textwidth,height=0.61\textheight,keepaspectratio]{fig_copper_gap.pdf}
  \vspace{-0.2em}

  {\small
    $NI/B=(1/\mu_0)L'+L/\mu_{\mathrm{Fe}}$;
    $\mu_0^{\mathrm{fit}}=8.8(7),\,8.32(34)\times10^{-7}\,\mathrm{T\,m/A}$.
  }
  \takeaway{The fits are linear, but both return $\mu_0$ about $30$--$34\%$ low.}
\end{frame}

\begin{frame}[plain]
  \centering
  \includegraphics[width=0.95\paperwidth,height=0.94\paperheight,keepaspectratio]{curie_scope_transition_frames.pdf}
\end{frame}

\begin{frame}{Curie: What The Raw Data Show}
  \centering
  \includegraphics[width=\textwidth,height=0.60\textheight,keepaspectratio]{curie_scope_transition_frames.pdf}
  \vspace{0.1em}

  {\small Series A: 141 stable loops retained over $180.8$--$273.3\,\mathrm K$; loops close as the sample warms.}
  \takeaway{The ferromagnetic--paramagnetic transition is visible before detailed fitting.}
\end{frame}

\begin{frame}{Curie: One Magnetization Number Per Loop}
  \small
  \begin{columns}[T,onlytextwidth]
    \begin{column}{0.32\textwidth}
      \textbf{Remanence}
      \[
        M_r=\frac12|M_+(0)-M_-(0)|
      \]
      \small Uses the loop opening at zero field.
    \end{column}
    \begin{column}{0.32\textwidth}
      \textbf{Common edge}
      \[
        M_{\mathrm{sat}}=\frac12|M_+(H_s)-M_-(-H_s)|
      \]
      \small Uses the same near-saturation field in every loop.
    \end{column}
    \begin{column}{0.32\textwidth}
      \textbf{Tail intercept}
      \[
        M_0=\frac12|b^{(+)}-b^{(-)}|
      \]
      \small Fits the high-field tails and subtracts the linear background.
    \end{column}
  \end{columns}
  \takeaway{Agreement locates the transition; disagreement measures analysis sensitivity.}
\end{frame}

\begin{frame}[plain]
  \centering
  \includegraphics[width=0.95\paperwidth,height=0.94\paperheight,keepaspectratio]{curie_method123_normalized.pdf}
\end{frame}

\begin{frame}{Curie Result: Apparent Transition}
  \centering
  \includegraphics[width=\textwidth,height=0.53\textheight,keepaspectratio]{curie_method123_normalized.pdf}
  \vspace{0.1em}

  {\small
    $M_r:216.4(8)\,\mathrm K$,
    $M_{\mathrm{sat}}:214.8(10)\,\mathrm K$,
    $M_0:199.2(17)\,\mathrm K$.
  }

  \vspace{0.1em}
  {\large $\Tapp=214(12)\,\mathrm K=-59(12)^\circ\mathrm C$}
  \takeaway{Finite-field half-height midpoint, not zero-field thermodynamic $\Tc$.}
\end{frame}

\begin{frame}[plain]
  \centering
  \includegraphics[width=0.95\paperwidth,height=0.94\paperheight,keepaspectratio]{curie_measured_drive_comparison.pdf}
\end{frame}

\begin{frame}{Curie: Field And Coverage Sensitivity}
  \centering
  \includegraphics[width=\textwidth,height=0.60\textheight,keepaspectratio]{curie_measured_drive_comparison.pdf}
  \vspace{0.05em}

  {\small
    Run means: A $214.0\,\mathrm K$, B $202.9\,\mathrm K$, C $236.5\,\mathrm K$;
    including C expands the envelope to $214(20)\,\mathrm K$.
  }
  \takeaway{Drive field and temperature coverage change the apparent midpoint.}
\end{frame}

\begin{frame}{What Limited The Measurements}
  \small
  \begin{columns}[T,onlytextwidth]
    \begin{column}{0.31\textwidth}
      \textbf{Copper-gap $\mu_0$}
      \begin{itemize}
        \item Effective gap thickness and plate contact.
        \item Flux leakage or wrong effective area.
        \item Integrator calibration.
        \item Incomplete saturation.
      \end{itemize}
    \end{column}
    \begin{column}{0.31\textwidth}
      \textbf{Curie midpoint}
      \begin{itemize}
        \item Finite-field rounding.
        \item Thermal lag and gradients.
        \item Linear background subtraction.
        \item Saturation-tail fit window.
      \end{itemize}
    \end{column}
    \begin{column}{0.31\textwidth}
      \textbf{Most useful repeats}
      \begin{itemize}
        \item Measure every copper plate directly.
        \item Calibrate the $B$ channel independently.
        \item Scan Curie more slowly.
        \item Ensure a cold plateau in every run.
      \end{itemize}
    \end{column}
  \end{columns}
  \takeaway{The next improvement is calibration and control, not more decimal places.}
\end{frame}

\begin{frame}{Conclusions}
  \large
  \begin{itemize}
    \item The iron toroid showed the expected soft-ferromagnet envelope and saturation, with $\mur^{\max}=2.18(8)\times10^3$.
    \item The copper-gap scans were linear, but returned $\mu_0$ values $30$--$34\%$ low; this identifies a scale systematic.
    \item The Monel loops collapsed near $\Tapp=214(12)\,\mathrm K=-59(12)^\circ\mathrm C$.
  \end{itemize}
  \takeaway{We observed both magnetic phenomena clearly; the honest final results are limited by systematic scale and finite-field effects.}
\end{frame}

\end{document}
